Collective climate action couples a public-goods dilemma with unequal vulnerability and contested redistribution.
Agents must decide how much of private wealth to transfer when peers can observe, gossip about, and sanction behavior, while formal rules may change through voting.
Classical theory treats reputation as a cooperation technology.
Visible giving raises standing.
Standing protects future partners and disciplines free-riding \pcite{milinski2002reputation,wedekind2000cooperation,nowak2006five}.
That directional claim now travels into societies of large language model (LLM) agents, where peer talk and scoring are easy to scaffold \pcite{park2023generative,piatti2024govsim,repunet2025}.

Unfortunately, LLM agents also exhibit calculated free-riding once payoff frames dominate deliberation \pcite{li2025spontaneous,freeriders2025}.
Communication and monitoring can therefore fail as repair devices even when average transfers look cooperative.
Climate asymmetry sharpens the test.
Developing agents may lack peer-enforcement tools.
Social stigma may land without a credible repair path.
Forced institutional assignment removes exit as an escape valve that earlier free-choice baselines allow.

In contrast, we treat social pressure as an event to be measured rather than as an assumed stabilizer.
We embed Meta Llama~3.1~8B agents in a repeated climate-risk economy with two parallel institutions, climatic shocks, a dual-use loss-and-damage fund, democratic whitelist voting, pairwise consistency scoring, aggregate reputation, and gossip bulletins.
Contribution intensity relative to wealth is the primary outcome.
We event-study intensity after bad-reputation marks, reputation drops, and gossip targeting.
We code post-event rationales for opportunistic versus repair motifs.

A single realization cannot separate structural behavior from stochastic path dependence.
Consequently, we replicate the full protocol under two independent seeds and report both runs with equal weight.
The multi-seed design yields a sharper claim than a one-shot narrative.
Across seeds, mean intensity after social marks falls for agents in the non-enforcing institution.
For agents in the enforcing institution, the gossip association flips from negative to positive.
Average cooperation levels differ across seeds, yet institution means remain matched within each seed.
Democracy expands subsidies in both runs, while punishment weakening is adopted in only one.

Our contributions are fourfold.
First, we provide a multi-seed event study of reputation and gossip under forced climate institutions.
Second, we separate robust withdrawal among non-enforcing agents from seed-sensitive responses among enforcing agents.
Third, we show that positive average cooperation coexists with limited path stability and sparse reputation-repair language.
Fourth, we document coverage without equalization in the loss-and-damage pool across both fiscal trajectories.
Together, these results challenge the default claim that social monitoring raises cooperation in LLM climate societies.
They also show why single-run narratives are insufficient for social-mechanism claims in stochastic LLM societies.
